\section{Notas de Clase}

\subsection{Adversaria y juegos}

En cuanto a la agenda es:
\begin{enumerate}
    \item Teoría de juegos: como juegos de cero-suma
    \item Decisiones optimas en juegos: algoritmo minimax, decisiones optimas en juegos multijugador, podada alfa-beta
    \item Busqueda de arbol alpha-beta heuristico: funciones de eval, cortando la busqueda, podada forward
\end{enumerate}

Las \textbf{Relajación de restricciones} es una técnica utilizada para simplificar problemas complejos de busqueda y optimización. Consiste en \textit{eliminar o debilitar} algunas de las reglas de un problema para encontrar rapidamente una \textbf{solución aproximada}.
Los pasos son los siguientes:
\begin{itemize}
    \item saber definir las variables y los dominios de valores
    \item saber definir las restricciones entre variables
    \item conocer la forma de resolución de este tipo de problemas y posibles heuristicas
\end{itemize}

Un ejemplo claro es el del sudoku, en este caso hay 8 reinas(8 variables) cada una tiene su posición $(x,y)$ cuyo diminio entero representa su posición en la fila/columna $1..8$
Cada variable peude tener un \textbf{dominio distinto}
Las \textbf{restricciones} las reinas no pueden compartir fila ni columna ni diagonal. En el sudoku los numeros no pueden repetirse en filas ni columnas ni en bloques de 3x3.
Pueden involucrar a tantas variables como queramos (binarias, terniarias, etc)

\subsection{Como aborda un CSP?}
Generamos un espacio de busqueda, y analizamos cada variable/mejorar la complejidad:
Hay un conjunto de \textbf{heuríasticas} con distintas aproximaciones $(tecnicas de inferencia)$ y durante la busqueda \textit{propagración y consistencia} y \textit{criterios de ordenacion y seleccion de variables y valores a instanciar}

\textbf{Util} cuando tenemos un problema que se puede modelar como combinación de variables y restricciones.
$x_1$, $x_2$ $\in$ [1..10] donde $x_1$ + $x_2$ > 10 y  $x_1$ - $x_2$ < 5.
Nos interessa encontrar cualquier solución que satisfaga todas las restricciones - satisfactibilidad. Proceso de reolución my caro. Elevada explosión combinatoria de heuristicas.


Es una extensión y hay relajación de restricciones. 
Juegos estocasticos
Los juegos clásicos nosotros estamos viendo es si las maquinas peuden jugar? Puedan tomar mejores decisiones que nosotros?
Los juegos son simulaciones, ambientes parcialemnte observables, tiene un numero de estados bastante grandes.


\subsection{Procesos de decisión markovianas}
\subsection{Tipos de juegos}
Tipo de objetivos (First person S)
Hay distintos tipos de ejes:
\begin{itemize}
    \item Deterministico o estocasticos
    \item uno, dos o mas jugadores
\end{itemize}

\subsubsection{Juegos de cero suma}
Agentes tienen utilidades opuestas (valores en outcomes)
Se puede pensar en un valoe unico que uno maximiza y el otro minimiza

\begin{definition}
    Tenemos una utilidad que \textbf{U=A,B}, donde A maximiza y B minimiza
\end{definition}

\textit{Cooperación implicita}
En los juegos generales los agentes tienen utilidades independientes, cooperación indiferencia competencia -> posibles

Agente toma decisiones segun el punto de desición, la utilidad es como la recompensa

\subsection{Algoritmo MiniMax}
Se aplica para juegos deterministicos y de cero suma. Hay un arbol de espacio de estados. Jugador alterna turnos. Computa el valor minimax de cada nodo (Mejor Utilidad alcanzable contra un jugador racional)

Recompenssas $\neq$ Busqueda

U = $f(n)$ $g=$ camino recorrido y $h=$ heuristica

\subsubsection{Desiciones optimas en Juegos}

Casos de cooperación implicita, donde hay multiples jugadores se toma en cuenta el costo por accion y considerar que hace al agente adversario

\subsection{Limites de Recursos}
Son juego srealisticos no puede buscar a hojas
La solución es Depht-limiteed search, lo cual busca solucionar para un esapcio limitado, reemplazar las utilidades terminales con una funcion de evaluacion para posiciones no terminales
Podada Forward, uso de estadisticas. Algoritmo de corte probabilistico y reduccion de movimientos tardios
